\begin{textsample}{POS Dim 2 – global\_south\_2025\_06 – Score 29.00 – global\_south\_2025\_06/126241208.txt}  \label{ex:f2_pos_242}
Before 1979, Iran was a linchpin of Western strategy in the Middle East.
Under Shah Mohammad Reza Pahlavi ( 1925 -- 79 ), Tehran was one of the United States ' closest \textbf{regional} partners and an unofficial Israeli \textbf{ally}.
The 1953 CIA-backed coup that restored the Shah entrenched this alliance.
Iran received extensive American military aid, acquired \textbf{nuclear} assistance, and became a vital \textbf{oil} supplier to Israel.
The Shah ' s regime even hosted an Israeli embassy and built an Iran -- Israel \textbf{oil} pipeline as part of Israel ' s"periphery doctrine"to break its \textbf{regional} isolation.
In 1979, the Islamic Revolution dismantled the pro-Western monarchy.
Ayatollah Khomeini seized power, branding the U.S. the"Great Satan"and Israel the"Little Satan."Diplomatic ties were abruptly severed : \textbf{Iranian} students stormed the U.S. embassy, taking hostages for 444 days, and Jewish and Israeli staff were expelled.
Iran ' s new constitution denounced Israel ' s legitimacy and cast itself as the champion of the Palestinian cause.
This marked a @ @ @ @ @ @ @ @ @ @ Iran ' s New Path Through the 1980s, Iran built an"\textbf{axis} of resistance"against Israel and its Western \textbf{allies}.
During the Iran -- Iraq War ( 1980 -- 88 ), Tehran relied on \textbf{proxies} like Hezbollah and Palestinian Islamic Jihad.
In 1984, the U.S. formally designated Iran a state sponsor of terrorism.
Tehran ' s Islamic \textbf{Revolutionary} \textbf{Guard} Corps ( IRGC ) -- notably its Quds Force -- nurtured \textbf{militias} across Lebanon, Syria, Iraq, Yemen, and Gaza.
Israel, viewing this \textbf{axis} as existentially threatening, engaged in covert action and frequent \textbf{air} \textbf{strikes} \textbf{targeting} \textbf{Iranian-backed} assets. \textbf{Nuclear} \textbf{Tensions} and Covert Warfare From the 2000s onwards, Iran accelerated its uranium enrichment program.
Israel and the U.S. responded with high-stakes covert actions, including the Stuxnet cyber \textbf{attack}.
The 2015 Joint Comprehensive Plan of Action ( JCPOA ) offered a temporary reset, curbing Iran ' s program in return for sanctions relief.
But this detente collapsed in 2018 when President Trump exited the deal, opted for"maximum pressure, @ @ @ @ @ @ @ @ @ @ enrichment to 60 per cent by 2020 -- just shy of weapons-grade -- prompting condemnation from Western \textbf{capitals}.
Rise in Open Conflict Post-2020 Since 2020, hostilities have \textbf{escalated} into almost open warfare.
The U.S. \textbf{assassination} of IRGC Quds Force \textbf{commander} Qassem Soleimani in 2020 triggered \textbf{Iranian} \textbf{missile} \textbf{strikes} on American bases in Iraq.
Covert warfare continued with the killings of \textbf{nuclear} scientists, mysterious sabotage at \textbf{Iranian} ports, and reciprocal \textbf{attacks}.
Heart of the Crisis : Israel -- Iran \textbf{Missile} Exchanges From late 2023 onward, the conflict became overt.
Iran publicly supported Hamas during the Gaza war, prompting Israeli \textbf{strikes} on Iranian-linked \textbf{targets} in Syria and Lebanon.
The \textbf{assassinations} of Hamas ' s political chief and Hezbollah ' s leader in 2024 triggered Iran ' s largest-ever direct \textbf{missile} \textbf{barrage} at Israel.
In mid-2025, Israel resumed \textbf{targeted} airstrikes within Iran, focusing on \textbf{nuclear} and military facilities.
In June 2025, the U.S. \textbf{launched} its own limited airstrikes on \textbf{Iranian} \textbf{nuclear} sites -- its first military assault directly on \textbf{Iranian} territory -- underscoring how deeply Iran @ @ @ @ @ @ @ @ @ @ \textbf{Axis} Emerges Tehran ' s animosity toward the West has pushed it toward Moscow.
By early 2024, Iran supplied Russia with advanced Fateh-110 \textbf{missiles}, ammunition, and kamikaze \textbf{drones} for Ukraine.
Russia, in return, supplied Iran with high-end Russian weaponry -- including Su-35 \textbf{jets}, Mi-28 helicopters, and S-400 systems -- and even shared satellite-launch resources.
The two countries cemented a long-term trade deal in December 2023.
Strategic Implications of the Alliance This \textbf{Iranian} -- Russian partnership represents a deeper strategic realignment.
For Iran, it provides valuable military cover against U.S. pressure ; for Russia, it bolsters its own anti-Western posture.
U.S. officials now see Tehran and Moscow as a"de facto united front,"spanning from Crimea to the Persian Gulf, complicating efforts to isolate Tehran.
Israeli leaders have warned that this integration strengthens Iran ' s \textbf{ballistic} \textbf{missile} \textbf{capabilities} and undercuts international deterrence.
While Iran ' s ties with India have remained cordial -- anchored in shared commercial and energy interests like the Chabahar port project -- its influence @ @ @ @ @ @ @ @ @ @ a wary relationship with Pakistan, especially given the sectarian dynamics along their \textbf{border}.
For Iran, decades of isolation and sanctions have confirmed its distrust of the West and driven it toward self-reliance and defiance.
Whether diplomacy, containment, or military pressure wins out remains to be seen -- but one thing is clear : Iran has rewritten its role in Middle Eastern geopolitics.

% matched lemmas: air, ally, assassination, attack, axis, ballistic, barrage, border, capability, capital, commander, drone, escalate, guard, iranian, iranian-backed, jet, launch, militia, missile, nuclear, oil, proxy, regional, revolutionary, strike, target, targeted, tension
\end{textsample}
